\section*{Capítulo \ref{ch:b3:representations} --- Representações de grupos
finitos}

\begin{solution}{exo:b3:representations:1}
(a) $\Z/n\Z$ é abeliano: todos os irredutíveis têm grau $1$
(\cref{prop:b3:representations:onedim}), isto é, são morfismos
$\chi \colon \Z/n\Z \to \C^\times$, determinados por $\chi(\bar 1)
= \omega$ com $\omega^n = 1$: os $n$ \omterm{def:b3:representations:character}{caracteres} $\chi_k(\bar
m) = \eu^{2\iu\pi km/n}$. Para $n = 4$ (classes $= $ elementos
$\bar0,\bar1,\bar2,\bar3$):
\[
\begin{array}{c|cccc}
& \bar0 & \bar1 & \bar2 & \bar3\\
\hline
\chi_0 & 1 & 1 & 1 & 1\\
\chi_1 & 1 & \iu & -1 & -\iu\\
\chi_2 & 1 & -1 & 1 & -1\\
\chi_3 & 1 & -\iu & -1 & \iu
\end{array}
\]
(b) Linhas: $\langle\chi_k,\chi_l\rangle = \frac14\sum_m
\eu^{2\iu\pi(l-k)m/4} = \delta_{kl}$ (soma geométrica); colunas do
mesmo modo. A matriz $(\eu^{2\iu\pi km/n})_{k,m}$ é a matriz da
\emph{transformada de Fourier discreta} --- o mesmo filtro de
raízes da unidade usado no capítulo de funções geradoras do
volume do segundo ano; a \omterm{thm:b3:representations:orthogonality}{ortogonalidade dos caracteres}
generaliza a fórmula de inversão da TFD.
\end{solution}

\begin{solution}{exo:b3:representations:2}
(a) A matriz de $\rho(g)$ na base $(e_x)$ é uma matriz de
permutação, de traço igual ao número de $x$ com $g\cdot x =
x$. Então
\[
\langle\mathbf 1, \chi\rangle = \frac1{\abs G}\sum_g
\abs{\operatorname{Fix}(g)} = \#\{\text{órbitas}\}
\]
pelo lema de contagem de Burnside (\cref{exo:b3:groups:5}) ---
equivalentemente, isso calcula a multiplicidade da \omterm{def:b3:representations:rep}{representação}
trivial, cujo espaço isotípico é o espaço dos vetores
$G$-invariantes, de dimensão igual ao número de \omterm{def:b3:groups:action}{órbitas} (um
indicador por \omterm{def:b3:groups:action}{órbita}).

(b) Como $\operatorname{Fix}_{X\times X}(g) =
\operatorname{Fix}_X(g)^2$, a parte (a) aplicada a $X \times X$
dá $\langle\chi,\chi\rangle = \frac1{\abs
G}\sum\abs{\operatorname{Fix}(g)}^2 = \#\{\text{órbitas em }
X\times X\}$ ($\chi$ é real). Escrevendo $\chi = \mathbf 1 +
\psi$: $\langle\mathbf1,\chi\rangle = 1$ (transitividade), de modo que
$\langle\psi,\psi\rangle = \langle\chi,\chi\rangle - 1$. A
ação sobre $X\times X$ tem a diagonal como uma \omterm{def:b3:groups:action}{órbita}; há
exatamente uma outra \omterm{def:b3:groups:action}{órbita} se, e somente se, $G$ é transitivo sobre os pares distintos:
$\langle\psi,\psi\rangle = 1$ se, e somente se, $2$-transitivo
(\cref{cor:b3:representations:multiplicity}).

(c) $S_n$ é $2$-transitivo sobre $\{1,\dots,n\}$ (leve qualquer
par distinto a qualquer outro): $\psi = \chi_{\mathrm{std}}$ é
\omterm{def:b3:representations:rep}{irredutível}.
\end{solution}

\begin{solution}{exo:b3:representations:3}
$S_3$: três classes, $\sum n_i^2 = 6 = 1 + 1 + 4$; os dois
\omterm{def:b3:representations:character}{caracteres} de grau $1$ são $\mathbf 1, \varepsilon$ ($[S_3 :
A_3] = 2$); a terceira linha $(2, a, b)$ decorre da
ortogonalidade de colunas com a coluna de $e$: $1 - 1 + 2a = 0$ e $1 + 1 +
2b = 0$: $a = 0$, $b = -1$ --- a tabela do~\cref{ex:b3:representations:s3}.

$S_4$, verificações (linhas): $\langle\chi_{\mathrm{std}},
\varepsilon\chi_{\mathrm{std}}\rangle = \frac1{24}(9 - 6 + 3 + 0
- 6) = 0$; $\langle\chi_2,\chi_2\rangle = \frac1{24}(4 + 0 + 12
+ 8 + 0) = 1$. Colunas: $e$ contra $(1\,2)$: $1 - 1 + 0 + 3 - 3
= 0$; $(1\,2)$ contra si mesma: $1 + 1 + 0 + 1 + 1 = 4 =
\abs{Z_{S_4}((1\,2))} = 24/6$.

\omterm{def:b3:representations:character}{Caráter} de permutação sobre $4$ pontos: $(4, 2, 0, 1, 0) = \mathbf
1 + \chi_{\mathrm{std}}$ (contagens de pontos fixos; subtraia a primeira
linha). Para $\chi_{\mathrm{std}}^2 = (9, 1, 1, 0, 1)$:
\[
\langle\mathbf1,\cdot\rangle = \tfrac{9 + 6 + 3 + 0 + 6}{24} =
1,\quad
\langle\varepsilon,\cdot\rangle = 0,\quad
\langle\chi_2,\cdot\rangle = \tfrac{18 + 6}{24} = 1,\quad
\langle\chi_{\mathrm{std}},\cdot\rangle = 1,\quad
\langle\varepsilon\chi_{\mathrm{std}},\cdot\rangle = 1:
\]
$\chi_{\mathrm{std}}^2 = \mathbf 1 + \chi_2 +
\chi_{\mathrm{std}} + \varepsilon\chi_{\mathrm{std}}$
(dimensões: $9 = 1 + 2 + 3 + 3$).
\end{solution}

\begin{solution}{exo:b3:representations:4}
Ambos os grupos têm cinco classes e padrão de graus $(1,1,1,1,2)$
(quatro de grau $1$ vindos da abelianização $\cong (\Z/2\Z)^2$,
e então $\sum n_i^2 = 8$). Ordenando as classes $e$, $z$ (a
involução central: $r^2$, resp.\ $-1$) e as três
classes de dois elementos:
\[
\begin{array}{c|ccccc}
& e & z & C_1 & C_2 & C_3\\
\hline
\chi^{(1)} & 1 & 1 & 1 & 1 & 1\\
\chi^{(2)} & 1 & 1 & 1 & -1 & -1\\
\chi^{(3)} & 1 & 1 & -1 & 1 & -1\\
\chi^{(4)} & 1 & 1 & -1 & -1 & 1\\
\chi^{(5)} & 2 & -2 & 0 & 0 & 0
\end{array}
\]
(a última linha vem da ortogonalidade de colunas). Tabelas idênticas para
$D_4$ e $Q_8$, que não são isomorfos
(\cref{pb:b3:groups:1}). A tabela \emph{de fato} captura: $\abs
G$, os tamanhos das classes, o \omterm{ex:b3:groups:actions}{centro} ($\{g : \abs{\chi_i(g)} = n_i\
\forall i\}$: ordem $2$ em ambos), a abelianização, todo o
reticulado dos \omterm{def:b3:groups:normal}{subgrupos normais} (núcleos e interseções,
\cref{exo:b3:representations:6}). Ela deixa escapar as ordens dos
elementos: $D_4$ tem cinco involuções, $Q_8$ tem uma --- de
modo que o tipo de isomorfismo é genuinamente mais fino que a
\omterm{def:b3:representations:table}{tabela de caracteres}.
\end{solution}

\begin{solution}{exo:b3:representations:5}
(a) Em $S_4$, o centralizador de $(1\,2\,3)$ tem ordem $24/8 =
3$: é $\langle(1\,2\,3)\rangle \subseteq A_4$. Logo
$Z_{A_4}((1\,2\,3))$ tem ordem $3$ e a $A_4$-classe tem $12/3
= 4$ elementos: os oito $3$-ciclos se repartem em duas $A_4$-classes
(representadas por $(1\,2\,3)$ e seu inverso).

(b) $A_4/V \cong \Z/3\Z$ dá três \omterm{def:b3:representations:character}{caracteres} de grau $1$
($\omega = \eu^{2\iu\pi/3}$; as classes de $3$-ciclos vão para
$\bar1, \bar2$); a restrição de $\chi_{\mathrm{std}}$
permanece \omterm{def:b3:representations:rep}{irredutível} ($\langle\chi,\chi\rangle = \frac1{12}(9 + 3
\cdot 1 + 0 + 0) = 1$):
\[
\begin{array}{c|cccc}
A_4 & e\,[1] & (1\,2)(3\,4)\,[3] & (1\,2\,3)\,[4] &
(1\,3\,2)\,[4]\\
\hline
\mathbf 1 & 1 & 1 & 1 & 1\\
\chi_\omega & 1 & 1 & \omega & \omega^2\\
\chi_{\bar\omega} & 1 & 1 & \omega^2 & \omega\\
\chi_3 & 3 & -1 & 0 & 0
\end{array}
\]
(c) Núcleos: $\ker\chi_\omega = \ker\chi_{\bar\omega} = V$;
$\ker\chi_3 = \{e\}$ (nenhuma outra entrada tem módulo $3$). Os
\omterm{def:b3:groups:normal}{subgrupos normais} são as interseções de núcleos
(\cref{exo:b3:representations:6}): $\{e\}$, $V$, $A_4$ --- em
particular, $A_4$ não tem \omterm{def:b3:groups:normal}{subgrupo normal} de ordem $2$ nem de índice
$2$.
\end{solution}

\begin{solution}{exo:b3:representations:6}
(a) Se $\chi(g) = \chi(e) = n$: a igualdade em
$\abs{\chi(g)} \leq n$ força $\rho(g) = \lambda\,\mathrm{id}$
(\cref{prop:b3:representations:charbasics}) com $n\lambda = n$:
$\rho(g) = \mathrm{id}$. A recíproca é clara.

(b) Seja $N \trianglelefteq G$. A \omterm{ex:b3:representations:examples}{representação regular} de
$G/N$ é fiel; decomponha-a em irredutíveis de $G/N$ e
levante-os a $G$ (\cref{prop:b3:representations:onedim}(b)):
\omterm{def:b3:representations:character}{caracteres} irredutíveis $\chi_{i_1}, \dots$ de $G$ cujos núcleos
contêm $N$ e cujo núcleo \emph{comum} é exatamente a
pré-imagem de $\{e\}$, isto é, $N$ (fidelidade no quociente).
Logo $N = \bigcap_j \ker\chi_{i_j}$.

(c) Se $G$ é simples: para um \omterm{def:b3:representations:rep}{irredutível} não trivial $\chi$,
$\ker\chi \trianglelefteq G$ não é $G$ (uma \omterm{def:b3:representations:rep}{representação
irredutível} trivial em todo $G$ é o \omterm{def:b3:representations:character}{caráter} trivial),
de modo que $\ker\chi = \{e\}$. Reciprocamente, suponha que todos os
núcleos não triviais sejam triviais, e seja $N \trianglelefteq G$ com $N \neq
G$. Na expressão de $N$ como interseção de núcleos em (b),
algum \omterm{def:b3:representations:character}{caráter} envolvido é não trivial (se todos fossem triviais, a
interseção seria $G$), e seu núcleo é $\{e\}$: $N =
\{e\}$. Logo os únicos \omterm{def:b3:groups:normal}{subgrupos normais} são $\{e\}$ e $G$.
\end{solution}

\begin{solution}{exo:b3:representations:7}
Ordem $8$ não abeliano: o número de \omterm{def:b3:representations:character}{caracteres} de grau $1$ é
$[G : D(G)]$, um divisor próprio de $8$ (não abeliano: $D(G) \neq
\{e\}$), e $\sum n_i^2 = 8$. Com $k$ deles iguais a um e os
graus restantes $\geq 2$: $8 - k \equiv 0$ com quadrados $\geq 4$, e $k
\mid 8$, $k < 8$. $k = 4$: um grau $2$ --- coerente. $k =
2$: sobram $6$, que não é soma de quadrados $\geq 4$. $k = 1$:
impossível, pois $k = [G : D(G)] \geq 2$ --- $G/D(G)$ é um
$2$-grupo abeliano não trivial, já que $G$ é um $2$-grupo com
$D(G) \neq G$ pela \omterm{def:b3:groups:derived}{solubilidade} dos $p$-grupos
(\cref{ex:b3:groups:solvableexamples}). Logo o padrão é
$(1,1,1,1,2)$. Fidelidade de $\chi_5$: os quatro \omterm{def:b3:representations:character}{caracteres} de grau $1$
contêm todos $D(G)$ em seus núcleos; se $\ker\chi_5
\supseteq N \neq \{e\}$ para algum normal minimal $N$ ($N
\subseteq D(G)$ ou não --- tome $N \subseteq \ker\chi_5$
não trivial), então $N$ estaria em todos os cinco núcleos, cuja
interseção é trivial (a \omterm{ex:b3:representations:examples}{representação regular} é
fiel): contradição. Ordem $p^3$ não abeliano: $Z(G)$ tem
ordem $p$ (ordem $p^2$ tornaria $G/Z$ cíclico e $G$ abeliano),
$G/Z(G)$, de ordem $p^2$, é abeliano, de modo que $D(G) \subseteq Z(G)$,
e $D(G) \ne \{e\}$: $D(G) = Z(G)$, dando $p^2$ \omterm{def:b3:representations:character}{caracteres} de
grau $1$. Os graus restantes satisfazem $\sum n_i^2 = p^3 -
p^2$ com cada $n_i > 1$ dividindo $\abs G$
(\cref{pb:b3:representations:1}, questão 8), logo $n_i \in
\{p\}$ ($n_i = p^2$ ultrapassaria: $p^4 > p^3 - p^2$): exatamente $p
- 1$ \omterm{def:b3:representations:character}{caracteres} de grau $p$.
\end{solution}

\begin{solution}{exo:b3:representations:8}
Defina, para \omterm{def:b3:representations:rep}{representações} $\rho, \sigma$ de $G, H$ em $V, W$,
a \omterm{def:b3:representations:rep}{representação} $\rho \boxtimes \sigma$ de $G \times H$ sobre
$V \otimes W$ --- concretamente, sobre matrizes: $(\rho\boxtimes
\sigma)(g,h)$ é o produto de Kronecker $\rho(g)\otimes
\sigma(h)$, cujo traço é
$\operatorname{tr}\rho(g)\operatorname{tr}\sigma(h) =
\chi(g)\psi(h)$ (o produto de Kronecker de matrizes $A \otimes B$
tem traço $\operatorname{tr}A\operatorname{tr}B$: sua diagonal
é $a_{kk}b_{ll}$). Logo $\chi\psi$ é um \omterm{def:b3:representations:character}{caráter}, e
\[
\langle\chi\psi, \chi'\psi'\rangle_{G\times H}
= \frac{1}{\abs G\abs H}\sum_{g,h}
\overline{\chi(g)\psi(h)}\,\chi'(g)\psi'(h)
= \langle\chi,\chi'\rangle_G\,\langle\psi,\psi'\rangle_H
= \delta_{\chi\chi'}\delta_{\psi\psi'}.
\]
Em particular $\langle\chi\psi,\chi\psi\rangle = 1$: cada
$\chi\psi$ é \omterm{def:b3:representations:rep}{irredutível}
(\cref{cor:b3:representations:multiplicity}). Esses são
$r_Gr_H$ \omterm{def:b3:representations:character}{caracteres} irredutíveis distintos; as classes de
$G\times H$ são os produtos de classes ($(g,h) \sim (g',h')$
componente a componente), de modo que há $r_Gr_H$ delas: a lista está
completa (\cref{thm:b3:representations:numberirr}). Para
$(\Z/2\Z)^2$: os quatro \omterm{def:b3:representations:character}{caracteres} de sinal
$(\pm1)\otimes(\pm1)$ --- a tabela do bloco superior esquerdo
do~\cref{exo:b3:representations:4}. Um grupo abeliano finito é
um produto de grupos cíclicos
(\cref{cor:b3:modules:abelian}); seus \omterm{def:b3:representations:character}{caracteres} irredutíveis são
produtos dos cíclicos (\cref{exo:b3:representations:1}):
todos de grau $1$.
\end{solution}

\begin{solution}{exo:b3:representations:9}
(a) $D(A_5) = A_5$ (simples não abeliano), de modo que o único \omterm{def:b3:representations:character}{caráter} de grau $1$
é $\mathbf 1$
(\cref{prop:b3:representations:onedim}). Precisamos de $n_2^2 + n_3^2
+ n_4^2 + n_5^2 = 59$ com cada $n_i \geq 2$; testando os
quadrados $4, 9, 16, 25, 36, 49$: o único multiconjunto que funciona é
$\{9, 9, 16, 25\}$: com maior quadrado $49$, o resto $10$
não é soma de três quadrados $\geq 4$; com $36$, o resto
$23$ também não é ($16 + 4 + 4 = 24$, $9 + 9 + 4 = 22$); com
maior $25$, verifica-se que $25 + 16 + 9 + 9 = 59$ funciona e que as variantes $25 +
25$, $25 + 16 + 16$, $25 + 16 + 4$ falham; com maior
$16$: $16\cdot3 = 48 < 59 - 4$. Graus: $1, 3, 3, 4, 5$.

(b) Permutação sobre $5$ pontos: pontos fixos $(5, 1, 2, 0, 0)$,
de modo que $\chi_4 = (4, 0, 1, -1, -1)$ com
$\langle\chi_4,\chi_4\rangle = \frac{16 + 0 + 20 + 12 + 12}{60}
= 1$: \omterm{def:b3:representations:rep}{irredutível}. Ação sobre os seis $5$-subgrupos de Sylow: uma
involução fixa exatamente $2$ (os \omterm{ex:b3:groups:actions}{normalizadores} são diedrais de
ordem $10$, cada um contendo $5$ involuções: $30$ incidências
para $15$ involuções), um elemento de ordem $3$ fixa $0$ (não há ordem $3$ em
$D_5$), um elemento de ordem $5$ fixa exatamente $1$ (ele está num único
Sylow): \omterm{def:b3:representations:character}{caráter} de permutação $(6, 2, 0, 1, 1)$, e $\chi_5 =
(5, 1, -1, 0, 0)$ com norma $\frac{25 + 15 + 20 + 0 + 0}{60} =
1$: \omterm{def:b3:representations:rep}{irredutível}. Restam duas linhas $(3, a, b, c, d)$, $(3, a', b', c',
d')$. Normas de colunas
(\cref{cor:b3:representations:column}): na classe de
$(1\,2)(3\,4)$, $\abs{Z} = 4$: $1 + 0 + 1 + a^2 + a'^2 = 4$;
coluna contra $e$: $1\cdot1 + 4\cdot 0 + 5\cdot1 + 3(a + a') =
0$: $a + a' = -2$, $a^2 + a'^2 = 2$: $a = a' = -1$. Nos
$3$-ciclos, $\abs Z = 3$: $1 + 1 + 1 + b^2 + b'^2 = 3$: $b = b'
= 0$. Em cada $5$-classe, $\abs Z = 5$: a coluna contra $e$
dá $1\cdot 1 + 4\cdot(-1) + 5\cdot 0 + 3(c + c') = 0$, de modo que
$c + c' = 1$; e
$c^2 + c'^2 = 3$: $\{c, c'\} = \bigl\{\frac{1+\sqrt5}2,
\frac{1-\sqrt5}2\bigr\}$ --- a razão áurea e sua
conjugada; a segunda $5$-classe carrega os valores trocados (as
duas linhas devem ser ortogonais).

(c) Na tabela terminada, nenhuma entrada de uma linha não trivial iguala
seu grau fora da primeira coluna: todo núcleo $\{g :
\chi_i(g) = n_i\}$ é trivial. Pelo~\cref{exo:b3:representations:6}(c), $A_5$ é simples.
\end{solution}

\begin{solution}{exo:b3:representations:10}
$\rho(z)$ comuta com todo $\rho(g)$ ($z$ é central), isto é,
$\rho(z) \in \operatorname{End}_G(V) = \C\,\mathrm{id}$ (Schur):
$\rho(z) = \lambda_z\,\mathrm{id}$, e $z \mapsto \lambda_z$ é
multiplicativa: um morfismo $Z(G) \to \C^\times$. Então
$\chi(z) = n\lambda_z$ com $\abs{\lambda_z} = 1$ (raiz da
unidade): $\abs{\chi(z)} = n$. Se $\rho$ é fiel, $\lambda$ é
injetiva em $Z(G)$ ($\rho(z) = \mathrm{id} \iff \lambda_z =
1$), de modo que $Z(G)$ mergulha em $\C^\times$; um subgrupo finito do grupo
multiplicativo de um corpo é cíclico
(\cref{thm:b3:galois:cyclic}).
\end{solution}

\begin{solution}{exo:b3:representations:11}
(a) Equivariância: $\rho(h)p_i\rho(h)^{-1}$ reindexa a soma
($g \mapsto hgh^{-1}$, e $\chi_i$ é uma \omterm{def:b3:representations:character}{função de classe}):
$p_i$ comuta com a ação. Sobre um \omterm{def:b3:representations:rep}{irredutível} $W$ de
\omterm{def:b3:representations:character}{caráter} $\chi_j$, Schur torna a restrição um escalar
$\lambda\,\mathrm{id}_W$; tomando traços,
\[
\lambda\,n_j = \frac{n_i}{\abs G}\sum_g
\overline{\chi_i(g)}\,\chi_j(g) = n_i\,\langle\chi_i,
\chi_j\rangle = n_i\,\delta_{ij}
\]
(primeira relação de ortogonalidade): $\lambda = \delta_{ij}$.

(b) Decomponha $V$ em irredutíveis (Maschke): $p_i$ age como a
identidade nas parcelas de \omterm{def:b3:representations:character}{caráter} $\chi_i$ e como $0$ em
todas as outras, de modo que $p_i$ é a projeção sobre a soma delas, $V_i$,
ao longo da soma das restantes; a imagem $V_i$ não depende da
decomposição escolhida (é o conjunto dos vetores fixados por
$p_i$, definido sem escolhas). $\sum_ip_i$ age como a identidade
em toda parcela \omterm{def:b3:representations:rep}{irredutível}: é $\mathrm{id}_V$. O
desmembramento mais fino de $V_i \cong W_i^{\oplus m_i}$ envolve
escolher uma base de $\operatorname{Hom}_G(W_i, V)$: canônico
ele não é.

(c) Para $\varepsilon$ (grau $1$): $p_\varepsilon =
\frac1{6}\sum_{g}\varepsilon(g)\,\rho(g)$, isto é, na álgebra
do grupo
\[
p_\varepsilon = \tfrac16\bigl(e - (1\,2) - (1\,3) - (2\,3)
+ (1\,2\,3) + (1\,3\,2)\bigr) .
\]
Elevando ao quadrado: o coeficiente de $g$ em $p_\varepsilon^2$ é
$\frac1{36}\sum_{h}\varepsilon(h)\varepsilon(h^{-1}g) =
\frac1{36}\,\varepsilon(g)\sum_h\varepsilon(h)^2 =
\frac{6}{36}\varepsilon(g)$: $p_\varepsilon^2 =
p_\varepsilon$. (Sua imagem na \omterm{ex:b3:representations:examples}{representação regular} é
a reta gerada por $\sum_g\varepsilon(g)e_g$: a \omterm{def:b3:representations:rep}{representação}
sinal aparece com multiplicidade $1$, como a teoria geral
exige.)
\end{solution}

\begin{solution}{exo:b3:representations:12}
(a) Ordene as classes $e$ [1], transposições [6],
$3$-ciclos [8], $4$-ciclos [6], transposições duplas [3].
O segundo \omterm{def:b3:representations:character}{caráter} linear é $\chi_2 = \varepsilon$, com
valores $1, -1, 1, -1, 1$. Para o \omterm{def:b3:representations:character}{caráter} de grau $2$
$\chi_3$, a ortogonalidade de cada coluna com a
coluna da identidade ($\sum_in_i\chi_i(g) = 0$ para $g \neq e$)
dá, nas transposições: $1 - 1 + 2\chi_3 + 3(\chi_4 +
\chi_5) = 0$; o truque do sinal $\chi_5 = \varepsilon\chi_4$
(um \omterm{def:b3:representations:character}{caráter} de grau $3$ vezes um linear é de novo
\omterm{def:b3:representations:rep}{irredutível} --- mesma norma) faz $\chi_4 + \chi_5$ se anular nas
classes ímpares: $\chi_3 = 0$ ali. Nos $3$-ciclos: $1 + 1 +
2\chi_3(c_3) + 3(\chi_4 + \chi_5)(c_3) = 0$ com
$\chi_5 = \chi_4$ nas classes pares; a coluna de $c_3$ consigo
mesma dá informação sobre $\abs{\chi_3}^2$; resolvendo o pequeno sistema
(use também a ortogonalidade de linhas de $\chi_3$ com $\mathbf 1$ e
$\varepsilon$): $\chi_3 = (2, 0, -1, 0, 2)$, e então $\chi_4 =
(3, 1, 0, -1, -1)$ e $\chi_5 = \varepsilon\chi_4 = (3, -1,
0, 1, -1)$. A tabela completa:
\begin{center}
\begin{tabular}{c|ccccc}
 & $e$ & $6\,t$ & $8\,c_3$ & $6\,c_4$ & $3\,v$\\
\hline
$\chi_1$ & $1$ & $1$ & $1$ & $1$ & $1$\\
$\chi_2$ & $1$ & $-1$ & $1$ & $-1$ & $1$\\
$\chi_3$ & $2$ & $0$ & $-1$ & $0$ & $2$\\
$\chi_4$ & $3$ & $1$ & $0$ & $-1$ & $-1$\\
$\chi_5$ & $3$ & $-1$ & $0$ & $1$ & $-1$\\
\end{tabular}
\end{center}
(Todas as linhas têm norma $1$; todas as colunas são ortogonais:
as verificações passam.)

(b) Os dados de classes $1, 6, 8, 6, 3$ com esses graus são os
de $S_4$ (tipos de ciclo $e$, $2$, $3$, $4$, $2{+}2$).
Núcleos: $\ker\chi_2 = \{g : \varepsilon(g) = 1\} = A_4$
(classes $e, c_3, v$: $1 + 8 + 3 = 12$, índice $2$);
$\ker\chi_3 = \{g : \chi_3(g) = 2\}$ = classes $e, v$: o
grupo de Klein $V_4$, de ordem $4$, normal. $\chi_4, \chi_5$
são fiéis ($\chi_i(g) = n_i$ só em $e$). Interseções
de núcleos: $\{e\}$, $V_4$, $A_4$, $G$ --- pelo~\cref{exo:b3:representations:6}(b), esses são \emph{todos} os
\omterm{def:b3:groups:normal}{subgrupos normais} de $S_4$.

(c) Os \omterm{def:b3:representations:character}{caracteres} com $V_4 \subseteq \ker$ são $\chi_1,
\chi_2, \chi_3$: eles se fatoram através de $G/V_4$, um grupo de
ordem $6$ que possui graus irredutíveis $1, 1, 2$ --- a
tabela de $S_3$. Como a tabela do quociente \emph{é} um
invariante completo entre os dois grupos de ordem $6$
($\Z/6\Z$ teria seis \omterm{def:b3:representations:character}{caracteres} lineares), $G/V_4 \cong
S_3$: o quociente é visível dentro da tabela como o bloco
de linhas que contêm $V_4$ em seu núcleo.
\end{solution}

\begin{solution}{pb:b3:representations:1}
\textbf{1.} Se $\alpha^n + c_{n-1}\alpha^{n-1} + \dots + c_0 =
0$ ($c_i \in \Z$), então $\alpha^n \in \Z\text{-gerado}(1, \dots,
\alpha^{n-1})$, e, indutivamente, toda potência o é: $\Z[\alpha]$ é
gerado por $1, \alpha, \dots, \alpha^{n-1}$. Reciprocamente, seja
$\Z[\alpha] = \Z g_1 + \dots + \Z g_m$. Escreva $\alpha g_i =
\sum_j m_{ij}g_j$ com $M = (m_{ij}) \in M_m(\Z)$: o vetor
$g = (g_i)$ satisfaz $(\alpha I - M)g = 0$; multiplicando pela
matriz adjunta, $\det(\alpha I - M)\,g_i = 0$ para todo $i$,
e, como $1 \in \Z[\alpha]$ é uma $\Z$-combinação dos
$g_i$, $\det(\alpha I - M) = 0$: $\alpha$ é raiz do
mônico $\det(XI - M) \in \Z[X]$.

\textbf{2.} Se $\Z[\alpha]$ é gerado pelas potências de $\alpha$ até
$n-1$ e $\Z[\beta]$ pelas de $\beta$ até $m - 1$, então
$\Z[\alpha,\beta]$ é gerado pelos $nm$ produtos
$\alpha^i\beta^j$ (reduza qualquer monômio). Os subanéis
$\Z[\alpha + \beta]$ e $\Z[\alpha\beta]$ são $\Z$-submódulos
do $\Z$-módulo \omterm{def:b3:modules:free}{finitamente gerado} $\Z[\alpha, \beta]$,
logo \omterm{def:b3:modules:free}{finitamente gerados} ($\Z[\alpha,\beta]$, gerado por
$nm$ elementos, é imagem de $\Z^{nm}$; um submódulo se levanta
a um submódulo de $\Z^{nm}$, livre de posto $\leq nm$ pelo~\cref{thm:b3:modules:submodule}, e sua imagem gera). Pela
questão 1, $\alpha + \beta$ e $\alpha\beta$ são inteiros
\omterm{def:b3:galois:algebraic}{algébricos}.

\textbf{3.} Se $\frac pq$ (fração \omterm{def:b3:representations:rep}{irredutível}) é raiz de um polinômio mônico
com coeficientes inteiros de grau $n$, o teorema das raízes racionais
(elimine os denominadores: $p^n = -q(\cdots)$) dá $q \mid p^n$,
de modo que $q = \pm1$.

\textbf{4.} $\chi(g)$ é uma soma de raízes da unidade
(\cref{prop:b3:representations:charbasics}), cada uma um inteiro
\omterm{def:b3:galois:algebraic}{algébrico} (raiz de $X^N - 1$); conclua pela questão 2.

\textbf{5.} Para $h \in G$: $\rho(h)S_C\rho(h)^{-1} = \sum_{g\in
C}\rho(hgh^{-1}) = S_C$ ($C$ é uma classe). Por Schur, $S_C =
\omega_C\,\mathrm{id}$; tomando traços, $\abs C\,\chi(g_C) =
\omega_C\, n$.

\textbf{6.} $S_CS_{C'} = \sum_{x \in C, y \in C'}\rho(xy) =
\sum_{z \in G} a(z)\rho(z)$ com $a(z) = \#\{(x,y) \in C\times
C' : xy = z\}$. A conjugação por $h$ estabelece uma bijeção entre as soluções para
$z$ e as para $hzh^{-1}$: $a$ é uma \omterm{def:b3:representations:character}{função de classe} com
valores em $\N$, de modo que $S_CS_{C'} = \sum_{C''}a_{CC'C''}S_{C''}$.
Substituindo $S_C = \omega_C\,\mathrm{id}$ em toda parte e
identificando os escalares:
$\omega_C\omega_{C'} = \sum_{C''}a_{CC'C''}\,\omega_{C''}$.

\textbf{7.} Seja $M$ o $\Z$-módulo gerado por $1$ e por todos os
produtos $\omega_{C_1}\cdots\omega_{C_k}$; pela questão 6, todo
tal produto se reduz a uma $\Z$-combinação de $1$ e dos
$\omega_C$: $M$ é \omterm{def:b3:modules:free}{finitamente gerado}, e $\omega_C M
\subseteq M$ para cada $C$. Em particular, $\Z[\omega_C]
\subseteq M$ é \omterm{def:b3:modules:free}{finitamente gerado} (submódulo, como na questão
2), e a questão 1 torna $\omega_C$ um inteiro \omterm{def:b3:galois:algebraic}{algébrico}.

\textbf{8.} Para um $\chi$ \omterm{def:b3:representations:rep}{irredutível} de grau $n$:
\[
\frac{\abs G}{n} = \frac{\abs G}{n}\langle\chi,\chi\rangle
= \frac1n \sum_{g}\chi(g)\overline{\chi(g)}
= \sum_{C}\frac{\abs C\,\chi(g_C)}{n}\,\overline{\chi(g_C)}
= \sum_C \omega_C\,\overline{\chi(g_C)} .
\]
Cada $\overline{\chi(g_C)} = \chi(g_C^{-1})$ é um inteiro
\omterm{def:b3:galois:algebraic}{algébrico} (questão 4), de modo que o membro direito é um (questões 2 e
7); ele é racional, logo inteiro (questão 3): $n \mid \abs
G$.

\textbf{9.} Bézout: $u\abs C + vn = 1$ com $u, v \in \Z$.
Então
\[
\frac{\chi(g_C)}{n} = u\,\frac{\abs C\,\chi(g_C)}{n} +
v\,\chi(g_C) = u\,\omega_C + v\,\chi(g_C),
\]
um inteiro \omterm{def:b3:galois:algebraic}{algébrico}.

\textbf{10.} Suponha $0 < \abs{\chi(g_C)} < n$ e seja $\alpha
= \chi(g_C)/n$. Todos os valores $\chi(g)$ estão em $\Q(\zeta_N)$, $N =
\abs G$ (somas de raízes $N$-ésimas da unidade). Para $\sigma \in
\operatorname{Gal}(\Q(\zeta_N)/\Q)$: $\sigma$ leva raízes da
unidade em raízes da unidade ($\sigma(\zeta^k) = \zeta^{ak}$,
\cref{thm:b3:galois:cyclotomicirred}), de modo que $\sigma(\chi(g_C))$
é de novo uma soma de $n$ raízes da unidade: $\abs{\sigma(\alpha)}
\leq 1$; além disso, $\sigma(\alpha)$ é um inteiro \omterm{def:b3:galois:algebraic}{algébrico} (tem
o mesmo \omterm{def:b3:galois:algebraic}{polinômio minimal} que $\alpha$). O produto $P =
\prod_\sigma \sigma(\alpha)$ é fixado por todo o \omterm{def:b3:galois:galois}{grupo de
Galois}, logo é racional (\cref{thm:b3:galois:fundamental}(1)),
e é um inteiro \omterm{def:b3:galois:algebraic}{algébrico} com
\[
0 < \abs P \leq \abs\alpha < 1
\]
(nenhum fator se anula: $\sigma(\alpha) = 0$ forçaria $\alpha =
0$). Isso contradiz a questão 3. Logo $\chi(g_C) = 0$ ou
$\abs{\chi(g_C)} = n$, e neste último caso $\rho(g_C)$ é
escalar (\cref{prop:b3:representations:charbasics}).

\textbf{11.} Ortogonalidade de colunas ($C \neq \{e\}$):
$\sum_i \chi_i(e)\overline{\chi_i(g_C)} = 0$, isto é, $1 +
\sum_{i \geq 2} n_i\overline{\chi_i(g_C)} = 0$. Se todo
$\chi_i$ não trivial com $p \nmid n_i$ se anulasse em $g_C$, então,
agrupando os demais por seu fator $p$:
\[
-\frac1p = \sum_{i \geq 2,\ p \mid n_i}
\frac{n_i}{p}\,\overline{\chi_i(g_C)},
\]
um inteiro \omterm{def:b3:galois:algebraic}{algébrico} --- contradizendo a questão 3. Logo algum
$\chi_i$ não trivial tem $p \nmid n_i$ e $\chi_i(g_C) \neq 0$;
como $\abs C = p^k$, $\gcd(\abs C, n_i) = 1$, e a questão 10
torna $\rho_i(g_C)$ escalar. Agora, $G$ simples não abeliano:
$\chi_i$ é fiel (\cref{exo:b3:representations:6}(c)), e
$Z_i = \{g : \rho_i(g) \text{ escalar}\}$ é um \omterm{def:b3:groups:normal}{subgrupo normal}
(a pré-imagem por $\rho_i$ dos escalares, que formam um \omterm{def:b3:groups:normal}{subgrupo
normal} --- de fato central --- da imagem) que contém
$g_C \neq e$: $Z_i = G$. Então $\rho_i(G)$ é abeliana e
fiel, tornando $G$ abeliano: contradição. \emph{Nenhum \omterm{def:b3:groups:simple}{grupo simples}
não abeliano tem \omterm{ex:b3:groups:actions}{classe de conjugação} de tamanho potência de primo
$> 1$.}

\textbf{12.} Seja $G$ simples de ordem $p^aq^b$. Se $b = 0$
($G$ um $p$-grupo): $Z(G) \neq \{e\}$
(\cref{thm:b3:groups:pfixed}) é normal, logo $Z(G) = G$: abeliano.
Caso contrário, tome $Q$ um $q$-subgrupo de Sylow e $z \in Z(Q)
\setminus\{e\}$ (\cref{thm:b3:groups:pfixed} de novo). Então $Q
\subseteq Z_G(z)$, de modo que a classe de $z$ tem tamanho $[G : Z_G(z)]$
dividindo $[G : Q] = p^a$: uma potência de $p$. Se o tamanho é $1$,
$z \in Z(G)$: o \omterm{ex:b3:groups:actions}{centro} é um \omterm{def:b3:groups:normal}{subgrupo normal} não trivial, logo
$Z(G) = G$, abeliano. Se o tamanho é $p^k > 1$: a questão 11
o proíbe para $G$ simples não abeliano. De um jeito ou de outro, um \omterm{def:b3:groups:simple}{grupo simples}
de ordem $p^aq^b$ é abeliano ($\cong \Z/p\Z$).

\textbf{13.} Indução sobre $\abs G$ ($\abs G = 1$: \omterm{def:b3:groups:derived}{solúvel}). Se
$G$ é simples, a questão 12 o torna abeliano, logo \omterm{def:b3:groups:derived}{solúvel}.
Caso contrário, tome $N$ um \omterm{def:b3:groups:normal}{subgrupo normal} próprio não trivial: $\abs
N$ e $\abs{G/N}$ são de novo da forma $p^{a'}q^{b'}$ e
menores, de modo que $N$ e $G/N$ são \omterm{def:b3:groups:derived}{solúveis} por indução, e $G$ é
\omterm{def:b3:groups:derived}{solúvel} (\cref{prop:b3:groups:derived}). --- Com três
primos, o passo-chave falha: o índice de um \omterm{def:b3:groups:sylow}{subgrupo de Sylow} deixa de
ser potência de primo, de modo que a classe de um elemento central de um
Sylow não precisa ter tamanho potência de primo. E alguma falha é
inevitável: $A_5$, de ordem $2^2\cdot3\cdot5$, é simples e
não \omterm{def:b3:groups:derived}{solúvel}.

\textbf{14.} As classes e os tamanhos são os do~\cref{exo:b3:groups:11}(a). A
$S_5$-classe de $c$ tem tamanho
$24$; se ela permanecesse uma única $A_5$-classe, a contagem
órbita--estabilizador daria $\abs{Z_{A_5}(c)} = 60/24$, que não é inteiro ---
concretamente, $Z_{S_5}(c) = \langle c\rangle$ tem ordem $5$ e
está dentro de $A_5$, de modo que a $A_5$-classe de $c$ tem tamanho $60/5 =
12$: a $S_5$-classe se parte em duas ($c$ e $c^2$ só são conjugados
em $S_5$ por uma permutação ímpar).

\textbf{15.} Um único \omterm{def:b3:representations:character}{caráter} linear: uma \omterm{def:b3:representations:rep}{representação} de grau $1$
se fatora através de $G/D(G)$, e $D(A_5) = A_5$ ($A_5$ é simples
não abeliano, e $D(A_5)$ é normal e não trivial ---
$A_5$ não é abeliano). Logo $n_1 = 1$ e $n_2^2 + n_3^2 + n_4^2
+ n_5^2 = 59$ com cada $n_i \geq 2$. Quadrados disponíveis: $4, 9,
16, 25, 36, 49$. Uma soma de quatro deles igual a $59$: o
maior deve ser $25$ ($36 + 4 + 4 + 9 = 53 < 59$ não se
ajusta: $36 + 16 + 4 + 4 = 60$, $36 + 9 + 9 + 4 = 58$, $36 +
16 + 9 + 4 = 65$ --- nenhuma combinação com $36$ ou $49$ funciona),
e $59 - 25 = 34 = 16 + 9 + 9$ (única possibilidade: $16 + 16 + 4 =
36$, $25 + 9 + 4 = 38$, $25 + 4 + 4 = 33$): graus $1, 3, 3,
4, 5$.

\textbf{16.} $\langle\pi, \mathbf 1\rangle = \frac1{60}(5 +
15\cdot1 + 20\cdot2 + 0 + 0) = 1$ (uma \omterm{def:b3:groups:action}{órbita} --- contagem de
Burnside), e $\langle\pi, \pi\rangle = \frac1{60}(25 + 15 + 80 +
0 + 0) = 2$: $\pi$ contém o \omterm{def:b3:representations:character}{caráter} trivial uma vez, e seu
outro constituinte é um único \omterm{def:b3:representations:rep}{irredutível}. Logo $\chi_4 = \pi
- \mathbf 1$ é \omterm{def:b3:representations:rep}{irredutível}, de grau $4$, com valores $4, 0,
1, -1, -1$.

\textbf{17.} Sobre os pares, um elemento fixa $\{i,j\}$ se, e somente se, fixa
ou troca $i, j$: as contagens são $10$ ($e$), $2$ ($t =
(1\,2)(3\,4)$ fixa $\{1,2\}, \{3,4\}$), $1$ ($(1\,2\,3)$
fixa $\{4,5\}$), $0, 0$. Então $\langle\pi_{10},
\pi_{10}\rangle = \frac1{60}(100 + 15\cdot4 + 20\cdot1) = 3$:
três constituintes irredutíveis, cada um uma vez. E
$\langle\pi_{10}, \mathbf 1\rangle = \frac1{60}(10 + 30 + 20)
= 1$, $\langle\pi_{10}, \chi_4\rangle = \frac1{60}(40 + 0 +
20\cdot1\cdot1 + 0 + 0) = 1$: logo $\pi_{10} = \mathbf 1 + \chi_4
+ \chi$ com $\chi$ \omterm{def:b3:representations:rep}{irredutível} de grau $10 - 1 - 4 = 5$ e
valores $\chi_5 = \pi_{10} - \mathbf 1 - \chi_4 = (5, 1, -1, 0,
0)$.

\textbf{18.} Escreva $a, a'$ para os valores de $\chi_2, \chi_3$
em $t$, e $b, b'$ em $s = (1\,2\,3)$; os quatro são reais ($t$
e $s$ são conjugados a seus inversos). Coluna $t$ contra a
coluna $e$: $1 + 3a + 3a' + 4\cdot0 + 5\cdot1 = 0$, de modo que $a + a'
= -2$; coluna $t$ consigo mesma: $1 + a^2 + a'^2 + 0 + 1 =
\frac{60}{15} = 4$, de modo que $a^2 + a'^2 = 2$; logo $(a + a')^2 = 4
= 2 + 2aa'$ dá $aa' = 1$ e $a = a' = -1$. Coluna $s$
contra $e$: $1 + 3(b + b') + 4 - 5 = 0$ dá $b + b' = 0$;
coluna $s$ consigo mesma: $1 + b^2 + b'^2 + 1 + 1 = \frac{60}
{20} = 3$ dá $b = b' = 0$.

\textbf{19.} Coluna $c$ contra a coluna $e$: $1 + 3(x + y) +
4(-1) + 5\cdot0 = 0$, de modo que $x + y = 1$. Coluna $c$ contra a
coluna $c^2$ (classes distintas, ortogonais): $1 + xy + yx + 1
+ 0 = 0$, de modo que $xy = -1$. Assim $x, y$ resolvem $T^2 - T - 1 = 0$:
$\{x, y\} = \{\varphi, \bar\varphi\}$ com $\varphi = \frac{1 +
\sqrt5}2$. Coerência: $x^2 + y^2 = (x+y)^2 - 2xy = 3 =
\frac{60}{12} - 2$ --- em acordo com a identidade da coluna consigo mesma $1 +
x^2 + y^2 + 1 + 0 = 5$. A tabela:
\begin{center}
\begin{tabular}{c|ccccc}
 & $e$ & $15\,t$ & $20\,s$ & $12\,c$ & $12\,c^2$\\
\hline
$\chi_1$ & $1$ & $1$ & $1$ & $1$ & $1$\\
$\chi_2$ & $3$ & $-1$ & $0$ & $\varphi$ & $\bar\varphi$\\
$\chi_3$ & $3$ & $-1$ & $0$ & $\bar\varphi$ & $\varphi$\\
$\chi_4$ & $4$ & $0$ & $1$ & $-1$ & $-1$\\
$\chi_5$ & $5$ & $1$ & $-1$ & $0$ & $0$\\
\end{tabular}
\end{center}

\textbf{20.} $\norm{\chi_2}^2 = \frac1{60}\bigl(9 + 15\cdot1 +
0 + 12\varphi^2 + 12\bar\varphi^2\bigr) = \frac{9 + 15 +
36}{60} = 1$, usando $\varphi^2 + \bar\varphi^2 = (\varphi +
\bar\varphi)^2 - 2\varphi\bar\varphi = 1 + 2 = 3$. Analogamente,
$\langle\chi_2, \chi_3\rangle = \frac1{60}(9 + 15 + 0 +
12(2\varphi\bar\varphi)) = \frac{9 + 15 - 24}{60} = 0$.
Divisibilidade de \omterm{thm:b3:galois:finitefields}{Frobenius}: $1, 3, 3, 4, 5$ dividem todos $60$.
Os valores irracionais $\varphi, \bar\varphi$ são a diferença
visível em relação a $S_5$: em $S_5$, todo elemento é conjugado a
todos os geradores de seu grupo cíclico com o mesmo tipo de ciclo
--- em particular $c \sim c^2$ ---, forçando valores de \omterm{def:b3:representations:character}{caráter} racionais (de fato
inteiros); em $A_5$, a partição dos
ciclos de comprimento cinco abre a porta para $\Q(\sqrt5)$.

\textbf{21.} Um \omterm{def:b3:groups:normal}{subgrupo normal} $N$ é uma união de classes,
contém $e$, e $\abs N \mid 60$. As somas candidatas:
$1{+}15 = 16$, $1{+}20 = 21$, $1{+}12 = 13$, $1{+}24 = 25$,
$1{+}15{+}20 = 36$, $1{+}15{+}12 = 28$, $1{+}15{+}24 = 40$,
$1{+}20{+}12 = 33$, $1{+}20{+}24 = 45$, $1{+}12{+}12 = 25$,
$1{+}15{+}20{+}12 = 48$, $1{+}15{+}20{+}24 = 60 - 12 = \dots$,
listando todas as subsomas próprias que contêm $1$: nenhum dos $13, 16,
21, 25, 28, 33, 36, 40, 45, 48, 13{+}\dots$ divide $60$,
exceto $1$ ele próprio: $N = \{e\}$ ou $A_5$. Simplicidade, lida em
cinco números. E o critério da questão 11 também está visível: os
tamanhos de classe $15 = 3\cdot5$, $20 = 4\cdot5$, $12 = 4\cdot3$ são
todos compostos de dois primos --- nenhuma classe de tamanho potência de
primo, exatamente como o critério de Burnside exige de um \omterm{def:b3:groups:simple}{grupo simples}.

\textbf{22.} Uma rotação de $\R^3$ de ângulo $\theta$ tem
autovalores $1, \eu^{\iu\theta}, \eu^{-\iu\theta}$: traço $1 +
2\cos\theta$. Para $\theta = \frac{2\pi}5$: $2\cos\frac{2\pi}5
= \frac{\sqrt5 - 1}2$, de modo que o traço é $1 + \frac{\sqrt5-1}2 =
\frac{1+\sqrt5}2 = \varphi$. O elemento não trivial $\tau$ de
$\operatorname{Gal}(\Q(\sqrt5)/\Q)$, aplicado entrada a entrada a uma
\omterm{def:b3:representations:table}{tabela de caracteres}, leva \omterm{def:b3:representations:character}{caracteres} em \omterm{def:b3:representations:character}{caracteres} (ele comuta
com a álgebra que os define: $\tau\circ\chi$ é o \omterm{def:b3:representations:character}{caráter} da
\omterm{def:b3:representations:rep}{representação} obtida transportando as matrizes por
$\tau$ nas entradas; ou, abstratamente: as relações de ortogonalidade
são $\Q$-racionais, de modo que $\tau$ permuta suas soluções);
$\tau$ fixa $\chi_1, \chi_4, \chi_5$ (valores racionais) e
deve, portanto, trocar $\chi_2$ e $\chi_3$: o segundo
\omterm{def:b3:representations:character}{caráter} de grau $3$ carrega os valores conjugados, sem matriz
alguma calculada. Geometricamente, as duas \omterm{def:b3:representations:rep}{representações} são a
ação icosaédrica e sua composta com um automorfismo externo
de $A_5$ (conjugação por uma transposição), que
troca as duas classes de ciclos de comprimento cinco.

\textbf{23.} Para $z \in Z(G)$, $\rho(z)$ comuta com todo
$\rho(g)$, de modo que, pelo lema de Schur, $\rho(z) = \lambda\,
\mathrm{id}$; como $z$ tem ordem finita, $\lambda$ é uma raiz
da unidade, e $\abs{\chi(z)} = \abs\lambda\,n = n$. Então
\[
\abs G = \abs G\,\langle\chi, \chi\rangle
= \sum_{g \in G}\abs{\chi(g)}^2
\geq \sum_{z \in Z(G)}\abs{\chi(z)}^2
= \abs{Z(G)}\,n^2,
\]
isto é, $n^2 \leq [G : Z(G)]$. Um grupo não abeliano de ordem $8$
($D_4$ ou $Q_8$) tem cinco \omterm{ex:b3:groups:actions}{classes de conjugação}, logo cinco
graus irredutíveis com $\sum n_i^2 = 8$; a única maneira de
escrever $8$ como soma de cinco quadrados $\geq 1$ é $1 + 1 + 1 + 1
+ 4$: graus $1, 1, 1, 1, 2$. Ambos os grupos têm \omterm{ex:b3:groups:actions}{centro} de
ordem $2$, e o \omterm{def:b3:representations:character}{caráter} de grau $2$ atinge a igualdade:
$2^2 = 4 = [G : Z(G)]$ --- a cota é atingida. Para $A_5$,
$Z = \{e\}$ e a cota se lê $n^2 \leq 60$: satisfeita por
$1, 3, 3, 4, 5$ com folga ($25 \leq 60$), como tem de
ser, já que a igualdade forçaria (pela mesma cadeia) $\chi$ a
se anular fora do \omterm{ex:b3:groups:actions}{centro}.

\textbf{24.} $\rho(g)^m = \mathrm{id}$ para $m$ a ordem de
$g$, de modo que $\rho(g)$ é anulado por $X^m - 1$, cindido com
raízes simples sobre $\C$: diagonalizável, com autovalores
$\lambda_1, \dots, \lambda_n$ raízes da unidade, numa base de
autovetores $(e_i)$. Os produtos $e_ie_j$ ($i \leq j$) formam uma
base de autovetores de $\operatorname{Sym}^2V$ com autovalores
$\lambda_i\lambda_j$, e $e_i \wedge e_j$ ($i < j$) uma de
$\Lambda^2V$; como
\[
\sum_{i<j}\lambda_i\lambda_j
= \frac{(\sum_i\lambda_i)^2 - \sum_i\lambda_i^2}2
= \frac{\chi(g)^2 - \chi(g^2)}2,
\]
e a soma simétrica acrescenta $\sum_i\lambda_i^2$ em vez de
subtraí-lo, ambas as fórmulas seguem. Para $\chi_2 = (3, -1,
0, \varphi, \bar\varphi)$ nas classes $(e, (2,2)\text{-},
3\text{-ciclos}, c, c^2)$: elevar ao quadrado leva as transposições
duplas a $e$, os três-ciclos em três-ciclos, a
classe de $c$ sobre a de $c^2$ e reciprocamente ($c^{-1} \sim
c$ em $A_5$, de modo que $c^4 \sim c$). Logo $\chi_2(g^2)$ se lê $(3,
3, 0, \bar\varphi, \varphi)$ e, usando $\varphi^2 = \varphi
+ 1$, $\bar\varphi = 1 - \varphi$:
\[
\Lambda^2\chi_2 = (3, -1, 0, \varphi, \bar\varphi) = \chi_2,
\qquad
\operatorname{Sym}^2\chi_2 = (6, 2, 0, 1, 1) .
\]
Decompondo o segundo, com os tamanhos de classe $1, 15, 20, 12,
12$: $\langle\cdot, \mathbf 1\rangle = \frac1{60}(6 +
15\cdot2 + 0 + 12 + 12) = 1$; $\langle\cdot, \chi_5\rangle =
\frac1{60}(30 + 30) = 1$; $\langle\cdot, \chi_4\rangle =
\frac1{60}(24 - 12 - 12) = 0$; $\langle\cdot, \chi_2\rangle =
\frac1{60}(18 - 30 + 12(\varphi + \bar\varphi)) = 0$, e
do mesmo modo para $\chi_3$. Assim $\operatorname{Sym}^2\chi_2 =
\mathbf 1 + \chi_5$ (dimensões $6 = 1 + 5$) e
$\chi_2\otimes\chi_2 = \mathbf 1 + \chi_2 + \chi_5$
(dimensões $9 = 1 + 3 + 5$). Geometria: o isomorfismo
equivariante $\Lambda^2\R^3 \to \R^3$, $u \wedge v \mapsto u
\times v$, é exatamente $\Lambda^2\chi_2 = \chi_2$ para um
grupo de rotações; a parcela $\mathbf 1$ do quadrado
simétrico é a forma quadrática invariante $x^2 + y^2 + z^2$, e
$\chi_5$ vive no espaço de dimensão cinco dos tensores
simétricos de traço nulo (quadráticas harmônicas).

\textbf{25.} Coluna de $c$ contra a coluna de $c^2$:
\[
1\cdot1 + \varphi\bar\varphi + \bar\varphi\varphi +
(-1)(-1) + 0 = 1 - 1 - 1 + 1 + 0 = 0,
\]
como a ortogonalidade exige para classes distintas ($\varphi
\bar\varphi = -1$). Coluna de $c$ consigo mesma: $1 +
\varphi^2 + \bar\varphi^2 + 1 + 0 = 1 + 3 + 1 = 5 = 60/12 =
\abs{Z_{A_5}(c)}$. \omterm{def:b3:representations:character}{Caráter} regular $\sum_in_i\chi_i$ nas
quatro colunas diferentes da identidade:
\[
1 - 3 - 3 + 0 + 5 = 0, \qquad
1 + 0 + 0 + 4 - 5 = 0,
\]
nas transposições duplas e nos três-ciclos, e na classe
de $c$ (a de $c^2$ é sua conjugada \omterm{def:b3:galois:galois}{de Galois}):
\[
1 + 3\varphi + 3\bar\varphi - 4 + 0 = 1 + 3 - 4 = 0,
\]
usando $\varphi + \bar\varphi = 1$. A tabela passa por toda
auditoria: é a \omterm{def:b3:representations:table}{tabela de caracteres} de $A_5$.
\end{solution}
